\chapter{Inside the Processor}

\begin{khulasa}[: what will we learn in this chapter?]
\begin{itemize}
    \item What is inside the processor, and the job of each part.
    \item Registers and their types.
    \item The arithmetic and logic unit: how the processor adds two numbers with logic gates, and what flags are.
    \item The control unit: who decides what happens at each moment.
    \item The instruction set and the addressing modes.
    \item The steps of every instruction: fetch, decode, execute, write back.
\end{itemize}
\end{khulasa}

%================================================
\section{A look inside the processor}
%================================================

Let us open the processor and look inside. The figure shows a simple processor, but it has all the main
ideas of real processors.

\begin{center}
\begin{tikzpicture}[r/.style={reg, minimum width=1.5cm},
    bus/.style={line width=3pt, color=brandTeal!70!black},
    link/.style={thick, color=brandTeal!70!black},
    sig/.style={dashed, thick, color=orangeDeep, -{Stealth[length=2mm]}}]
    \draw[brandTeal, very thick, rounded corners=8pt, fill=brandTeal!4] (-8.4,-3.3) rectangle (2.2,3.4);
    \node[font=\bfseries, text=tealDeep, anchor=north east] at (2.1,3.35) {Processor};
    \draw[bus] (-7.6,0) -- (1.8,0);
    \node[font=\scriptsize, text=tealDeep] at (0.2,-0.3) {internal bus};
    % control unit
    \draw[orangeDeep, very thick, rounded corners=5pt, fill=brandOrange!12] (-8.1,0.5) rectangle (-4.9,3.0);
    \node[font=\small\bfseries, text=orangeDeep] at (-6.5,2.7) {control unit (UC)};
    \node[r] (ri) at (-6.5,2.05) {RI};
    \node[reg, font=\scriptsize, minimum width=1.25cm] (dec) at (-7.3,1.05) {decoder};
    \node[reg, font=\scriptsize, minimum width=1.25cm] (seq) at (-5.7,1.05) {sequencer};
    \draw[sahm] (ri) -- (dec);
    \draw[sahm] (dec) -- (seq);
    \draw[link] (ri.east) -- ++(0.35,0) |- (-4.55,0.35) -- (-4.55,0);
    \node[unit=brandPurple, font=\small] (clk) at (-5.8,3.95) {clock};
    \draw[sahm] (clk) -- (-5.8,3.0);
    % registers above the bus
    \node[r] (co) at (-2.6,2.3) {CO};
    \node[r] (ram) at (0.6,2.3) {RAM};
    \node[r] (rim) at (0.6,1.1) {RIM};
    \draw[sahmT] (co) -- node[above, font=\scriptsize, text=brandDark] {address} (ram);
    \node[font=\scriptsize, text=brandMuted, anchor=west] at (-2.45,1.65) {+1 per fetch};
    \draw[link] (co.south) ++(0,-0.15) -- (-2.6,0);
    \draw[link] (rim.south) -- (0.6,0);
    % ALU with two inputs, ACC under it
    \draw[brandGreen!70!black, very thick, fill=brandGreen!20]
        (-5.4,-0.8) -- (-3.6,-0.8) -- (-4.0,-1.9) -- (-5.0,-1.9) -- cycle;
    \node[font=\bfseries] at (-4.5,-1.35) {UAL};
    \node[r] (acc) at (-4.5,-2.75) {ACC};
    \node[r, font=\small] (fl) at (-2.3,-2.75) {flags};
    \draw[sahmT] (-4.5,-1.9) -- (acc.north) node[pos=0.5, right, font=\scriptsize, text=brandDark] {result};
    \draw[sahmT] (-3.8,-1.6) -| (fl.north);
    \draw[sahm] (-3.9,0) -- (-3.9,-0.8) node[pos=0.5, right, font=\scriptsize, text=brandDark] {operand 2};
    \draw[sahm] (acc.west) -- ++(-0.9,0) |- (-6.0,-0.45) -| (-5.1,-0.8);
    \node[font=\scriptsize] at (-6.75,-1.6) {operand 1};
    \draw[link] (acc.south) -- (-4.5,-3.2) -- (-1.0,-3.2) -- (-1.0,0);
    % control signals
    \draw[sig] (seq.south) -- (-5.7,0.15);
    \node[font=\scriptsize, text=orangeDeep] at (-7.3,0.25) {control signals};
    % memory
    \node[unit=brandPurple, minimum width=2cm, minimum height=4cm] (mem) at (4.4,1.2) {main\\memory};
    \draw[sahmO] (ram.east) -- node[above, font=\scriptsize, text=brandDark] {address bus} ([yshift=1.1cm]mem.west);
    \draw[sahmT, {Stealth[length=2.6mm]}-{Stealth[length=2.6mm]}] (rim.east) -- node[below, font=\scriptsize, text=brandDark] {data bus} ([yshift=-0.1cm]mem.west);
\end{tikzpicture}
\end{center}

\begin{center}
\small
\begin{tabularx}{\linewidth}{|>{\centering\arraybackslash}p{4cm}|R|}
    \hline
    \ras{Part} & \ras{Its job in one sentence} \\ \hline
    CO (program counter) & keeps the address of the \textbf{next instruction} \\ \hline
    RI (instruction register) & keeps the \textbf{current instruction} being run \\ \hline
    RAM (memory address register) & keeps the address we send to memory \\ \hline
    RIM (memory data register) & keeps the data coming from memory or going to it \\ \hline
    ACC (accumulator) & keeps the result of the calculation \\ \hline
    flags & keep information about the last result: is it zero? negative? \dots \\ \hline
    UAL & does the arithmetic and logic operations \\ \hline
    decoder and sequencer & understand the instruction and send the orders in the right order \\ \hline
\end{tabularx}
\end{center}

\begin{intabih}
    Do not mix up RAM the memory (\emph{Random Access Memory}) and RAM the register (\emph{Registre d'Adresse
    Mémoire}) used in French books. In English this register is called \textbf{MAR}, RIM is called \textbf{MDR},
    and CO is called \textbf{PC} (\emph{Program Counter}). In the 8086 it is called \textbf{IP}.
    In this course we keep the French short names (CO, RI, RAM, RIM, UAL, UC) because they are used in the syllabus.
\end{intabih}

%================================================
\section{Registers}
%================================================

\begin{tarif}
    A \term{register}{registre}: a very small memory box (8, 16, 32 or 64 bits) \textbf{inside the processor}.
    It is the fastest memory in the whole computer.
\end{tarif}

\begin{tashbih}
    When you solve a maths exercise, you do not walk to the library each time you need a number! You write
    the numbers you need on \textbf{scrap paper} in front of you. Registers are the processor's scrap paper,
    and memory is the library.
\end{tashbih}

\subsection{Why are registers important for speed?}
\begin{center}
\begin{tabular}{|c|c|c|}
    \hline
    \ras{Place} & \ras{Access time (about)} & \ras{If one cycle took one second\dots} \\ \hline
    register & less than 1 ns & one second \\ \hline
    main memory (RAM) & about 80 ns & about a minute and a half \\ \hline
    SSD disk & about 100 µs & more than a day \\ \hline
\end{tabular}
\end{center}
So we keep in registers the data we use often, and we go to memory as little as possible.

\subsection{Types of registers}
\begin{enumerate}
    \item \term{Data registers}{registres de données}: they keep the numbers we compute with (like AX, BX in the 8086).
    \item \term{Address registers}{registres d'adresse}: they keep memory addresses (like SI, DI, SP).
    \item \term{Control and status registers}{registres de contrôle}: they help the processor organise its work: CO, RI, and the flags register.
\end{enumerate}

\begin{mithal}[: registers at work]
    The processor is now running the instruction at address 11, ``add ACC and the content of box 21''.
    \begin{itemize}
        \item CO holds 12: the address of the \textbf{next} instruction, not the current one.
        \item RI holds the current instruction, so the control unit can read it during the whole execution.
        \item RAM holds 21: the address we are sending to memory now.
        \item RIM receives the content of box 21 coming from memory.
        \item ACC holds the first number, and will receive the result.
    \end{itemize}
    Each register has a clear job. This is why we call them special registers, unlike the general registers
    that the programmer uses freely.
\end{mithal}

\subsection{The flags register}
\begin{tarif}
    A \term{flag}{indicateur, drapeau}: one bit that the arithmetic unit sets to 1 or 0 after each
    operation, depending on the result. We use flags later to make decisions (for example: ``if the result is zero, jump'').
\end{tarif}

\begin{center}
\begin{tabularx}{\linewidth}{|>{\centering\arraybackslash}p{2.6cm}|R|R|}
    \hline
    \ras{Flag} & \ras{Becomes 1 when\dots} & \ras{Example on 8 bits} \\ \hline
    ZF (zero) & the result is \textbf{zero} & $5 - 5 = 0$ \\ \hline
    CF (carry) & a carry comes out of the last bit (the result is too big) & $200 + 100 = 300 > 255$ \\ \hline
    SF (sign) & the last bit of the result is 1 (negative number) & $3 - 7 = -4$ \\ \hline
    OF (overflow) & the result is wrong if we read the numbers as signed & $100 + 50 = 150 > 127$ \\ \hline
\end{tabularx}
\end{center}

\subsection{The difference between CF and OF}
The same byte can be read in two ways, and the processor does not know which one the programmer means:
\begin{itemize}
    \item \textbf{unsigned}: values from 0 to 255.
    \item \textbf{signed}: values from $-128$ to $+127$. The last bit (bit 7) is the sign: 1 means negative.
\end{itemize}
So the processor computes both flags. CF matters for unsigned numbers, and OF matters for signed numbers.

\begin{khatwa}[: $100 + 50$ on 8 bits]
    \begin{enumerate}
        \item $100 = \bin{0110\,0100}$ and $50 = \bin{0011\,0010}$.
        \item The sum: $\bin{1001\,0110} = 150$ (unsigned). No carry out of bit 7, so CF = 0: correct for unsigned numbers.
        \item But bit 7 is now 1! Reading it as signed, we see a negative number: $150 - 256 = -106$.
        \item We added two positive numbers and got a negative one: impossible, so the signed result is wrong, and OF = 1.
    \end{enumerate}
\end{khatwa}

\begin{intabih}
    An easy rule for OF: overflow happens when we add two numbers \textbf{with the same sign} and get a result
    \textbf{with the other sign}. Adding a positive and a negative number never overflows.
\end{intabih}

%================================================
\section{The arithmetic and logic unit (UAL)}
%================================================

\begin{tarif}
    The \term{arithmetic and logic unit}{unité arithmétique et logique, UAL} (ALU in English): the part
    that does arithmetic (add, subtract, compare\dots) and logic (AND, OR, XOR, NOT, shifts).
\end{tarif}

\begin{center}
\begin{tikzpicture}
    \draw[brandGreen!70!black, very thick, fill=brandGreen!20]
        (-2.2,1) -- (-0.4,1) -- (0,0.4) -- (0.4,1) -- (2.2,1) -- (1.2,-1) -- (-1.2,-1) -- cycle;
    \node[font=\bfseries\large] at (0,-0.2) {UAL};
    \draw[sahmT] (-1.3,2) node[above, font=\small] {operand A} -- (-1.3,1);
    \draw[sahmT] (1.3,2) node[above, font=\small] {operand B} -- (1.3,1);
    \draw[sahmO] (-3.6,0) node[left, font=\small] {operation (from the control unit)} -- (-1.7,0);
    \draw[sahmT] (0,-1) -- (0,-2) node[below, font=\small] {result};
    \draw[sahm] (1.6,-0.2) -- (3.4,-0.2) node[right, font=\small] {flags ZF, CF, SF, OF};
\end{tikzpicture}
\end{center}

\subsection{How do logic gates add? The half adder}
Let us add two bits $A$ and $B$. There are only four cases:

\begin{center}
\begin{tabular}{|c|c||c|c|}
    \hline
    \ras{$A$} & \ras{$B$} & \ras{sum $S$} & \ras{carry $C$} \\ \hline
    0 & 0 & 0 & 0 \\ \hline
    0 & 1 & 1 & 0 \\ \hline
    1 & 0 & 1 & 0 \\ \hline
    1 & 1 & 0 & \cell{orangeDeep}{\textbf{1}} \\ \hline
\end{tabular}
\end{center}

We see that $S$ is 1 when exactly one bit is 1: this is the \textbf{XOR} gate.
And $C$ is 1 when both bits are 1: this is the \textbf{AND} gate.

\begin{tarif}
    \term{Half adder}{demi-additionneur}: \quad $S = A \oplus B$ \quad and \quad $C = A \cdot B$
\end{tarif}

\begin{center}
\begin{tikzpicture}[gate/.style={draw=brandDark, fill=white, thick, minimum width=1.5cm,
        minimum height=0.9cm, font=\codefont\bfseries, rounded corners=2pt}]
    \node[gate] (x) at (0,0.8) {XOR};
    \node[gate] (a) at (0,-0.8) {AND};
    \node[font=\bfseries] (A) at (-3.2,0.4) {A};
    \node[font=\bfseries] (B) at (-3.2,-0.4) {B};
    \fill (-2,0.4) circle (2pt); \fill (-1.5,-0.4) circle (2pt);
    \draw[thick] (A) -- (-2,0.4) |- ([yshift=0.15cm]x.west);
    \draw[thick] (-2,0.4) |- ([yshift=0.15cm]a.west);
    \draw[thick] (B) -- (-1.5,-0.4) |- ([yshift=-0.15cm]x.west);
    \draw[thick] (-1.5,-0.4) |- ([yshift=-0.15cm]a.west);
    \draw[sahmT] (x.east) -- (2,0.8) node[right] {S};
    \draw[sahmO] (a.east) -- (2,-0.8) node[right] {C};
\end{tikzpicture}
\end{center}

\subsection{The full adder}
To add numbers with many bits, we must also add the \textbf{incoming carry} from the previous bit.
A full adder has three inputs: $A$, $B$ and the carry in $C_{in}$.

\begin{tarif}
    \term{Full adder}{additionneur complet}: \quad
    $S = A \oplus B \oplus C_{in}$ \quad and \quad
    $C_{out} = A \cdot B + C_{in} \cdot (A \oplus B)$
\end{tarif}

\begin{center}
\small
\begin{tabular}{|c|c|c||c|c|}
    \hline
    \ras{$A$} & \ras{$B$} & \ras{$C_{in}$} & \ras{$S$} & \ras{$C_{out}$} \\ \hline
    0&0&0&0&0 \\ \hline 0&0&1&1&0 \\ \hline 0&1&0&1&0 \\ \hline 0&1&1&0&1 \\ \hline
    1&0&0&1&0 \\ \hline 1&0&1&0&1 \\ \hline 1&1&0&0&1 \\ \hline 1&1&1&1&1 \\ \hline
\end{tabular}
\end{center}

\subsection{A 4-bit adder}
We put 4 full adders side by side: the carry out of each one goes into the next one, exactly as we do
when we add with pen and paper.

\begin{center}
\begin{tikzpicture}[fa/.style={draw=brandTeal, fill=brandTeal!15, very thick, minimum width=1.6cm,
        minimum height=1.1cm, font=\codefont\bfseries, rounded corners=3pt}]
    \foreach \i/\x in {0/0, 1/-3, 2/-6, 3/-9} {
        \node[fa] (f\i) at (\x,0) {FA};
        \draw[thick] ([xshift=-0.35cm]f\i.north) -- ++(0,0.6) node[above, font=\small] {A\i};
        \draw[thick] ([xshift=0.35cm]f\i.north) -- ++(0,0.6) node[above, font=\small] {B\i};
        \draw[sahmT] (f\i.south) -- ++(0,-0.6) node[below, font=\small] {S\i};
    }
    \draw[sahmO] (1.7,0) node[right, font=\small] {0} -- (f0.east);
    \draw[sahmO] (f0.west) -- (f1.east);
    \draw[sahmO] (f1.west) -- (f2.east);
    \draw[sahmO] (f2.west) -- (f3.east);
    \draw[sahmO] (f3.west) -- ++(-0.9,0) node[left, font=\small] {CF};
\end{tikzpicture}
\end{center}
Bit 0 (the smallest) is on the right, like when we write numbers. The carry travels from right to left.

\begin{khatwa}[: add $0111 + 0110$ (that is $7 + 6$) with the adder]
    \begin{center}
    \begin{tabular}{|c|c|c|c|c|c|}
        \hline
        \ras{Bit} & \ras{$A$} & \ras{$B$} & \ras{$C_{in}$} & \ras{$S$} & \ras{$C_{out}$} \\ \hline
        0 & 1 & 0 & 0 & 1 & 0 \\ \hline
        1 & 1 & 1 & 0 & 0 & 1 \\ \hline
        2 & 1 & 1 & 1 & 1 & 1 \\ \hline
        3 & 0 & 0 & 1 & 1 & 0 \\ \hline
    \end{tabular}
    \end{center}
    Read $S$ from bit 3 down to bit 0: $1101 = 13$, which is indeed $7 + 6$. The last carry is 0, so CF = 0.
\end{khatwa}

\subsection{And subtraction? The same adder!}
The processor does not need a special circuit to subtract. It uses the \term{two's complement}{complément à deux}:
$A - B = A + (-B)$, and we get $-B$ by flipping all the bits of $B$, then adding 1.

\begin{khatwa}[: $7 - 5$ on 4 bits]
    \begin{enumerate}
        \item $5 = 0101$. Flip the bits: $1010$. Add 1: $1011$. So $-5 = 1011$.
        \item Add: $0111 + 1011 = 1\,0010$. We drop the carry out, so $0010 = 2$ is left.
        \item The result: $7 - 5 = 2$. Correct! The same circuit added and subtracted.
    \end{enumerate}
\end{khatwa}

\subsection{Comparing by subtracting}
The processor does not need a special circuit to compare either: to compare $A$ and $B$ it computes
$A - B$ \textbf{without keeping the result}, then looks at the flags:
\begin{itemize}
    \item if ZF = 1 then $A = B$.
    \item if CF = 1 (a ``borrow'' happened) then $A < B$ (for unsigned numbers).
\end{itemize}

\subsection{A real test: how do the flags change?}
We ran the program below on the 8086. Do not worry if you do not understand all the instructions yet; we
study them in Chapter 3. Just look at the flags.

\asm{ch2/flags.asm}

\begin{center}
\small
\begin{tabular}{|c|l|c|c|c|c|c|}
    \hline
    \ras{Line} & \ras{Instruction} & \ras{AL} & \ras{CF} & \ras{ZF} & \ras{SF} & \ras{OF} \\ \hline
    4 & \code{add al, 100} ($200 + 100$) & $\hex{2C} = 44$ & \cell{orangeDeep}{\textbf{1}} & 0 & 0 & 0 \\ \hline
    5 & \code{mov al, 100} & $\hex{64} = 100$ & 1 & 0 & 0 & 0 \\ \hline
    6 & \code{add al, 50} ($100 + 50$) & $\hex{96} = 150$ & 0 & 0 & 1 & \cell{orangeDeep}{\textbf{1}} \\ \hline
    8 & \code{sub al, 5} ($5 - 5$) & $\hex{00} = 0$ & 0 & \cell{orangeDeep}{\textbf{1}} & 0 & 0 \\ \hline
    10 & \code{sub al, 7} ($3 - 7$) & $\hex{FC} = 252$ & \cell{orangeDeep}{\textbf{1}} & 0 & \cell{orangeDeep}{\textbf{1}} & 0 \\ \hline
\end{tabular}
\end{center}

\begin{khatwa}[: why did $200 + 100$ give 44?]
    \begin{enumerate}
        \item $200 + 100 = 300$, but AL holds only 8 bits, that is from 0 to 255.
        \item $300 - 256 = 44$: AL keeps 44, and the ``256'' that did not fit goes into the flag CF, so CF = 1.
    \end{enumerate}
\end{khatwa}

\begin{intabih}
    Look at line 5: the instruction \code{mov} \textbf{does not change the flags}, so CF stayed 1. Only
    arithmetic and logic instructions change the flags.
\end{intabih}

%================================================
\section{The control unit (UC)}
%================================================

\begin{tarif}
    The \term{control unit}{unité de commande, UC}: the part that reads the instructions and understands
    them, then sends \textbf{signals} to the other parts at each clock pulse (open this register, add, read
    memory\dots) in the right order.
\end{tarif}

\begin{tashbih}
    The control unit is the \textbf{orchestra conductor}: it does not play any instrument itself, but it tells
    each musician when to start and when to stop. Without it, everybody plays at the same time and it becomes noise.
\end{tashbih}

Its main parts:
\begin{itemize}
    \item the \textbf{instruction register} (RI): keeps the current instruction.
    \item the \term{decoder}{décodeur}: reads the operation code (opcode) and finds what is asked: add? move? jump?
    \item the \term{sequencer}{séquenceur}: sends the right signals, one after the other, to the rhythm of the clock.
\end{itemize}

\subsection{Micro-operations, pulse by pulse}
Each step of an instruction is split into very small steps called \term{micro-operations}{micro-opérations}.
At each clock pulse, the sequencer sends signals that open ``doors'' between registers. Here is how an instruction is fetched:

\begin{center}
\small
\begin{tabularx}{\linewidth}{|c|R|R|}
    \hline
    \ras{Pulse} & \ras{Micro-operation} & \ras{Signal sent by the sequencer} \\ \hline
    T1 & \nql{RAM}{CO} & ``copy CO into RAM'' \\ \hline
    T2 & \nql{RIM}{[RAM]}, \nql{CO}{CO + 1} & ``read memory'' and ``increase CO'' \\ \hline
    T3 & \nql{RI}{RIM} & ``copy RIM into RI'' \\ \hline
    T4\dots & depends on the instruction & the decoder chooses the next signals \\ \hline
\end{tabularx}
\end{center}
The fetch is the same for all instructions; what comes after depends on the opcode.

\subsection{Two kinds of control units}
\begin{center}
\begin{tabularx}{\linewidth}{|>{\centering\arraybackslash}p{3.2cm}|R|R|}
    \hline
    \ras{} & \ras{Hardwired (câblée)} & \ras{Microprogrammed (microprogrammée)} \\ \hline
    Idea & a fixed logic circuit makes the signals & each instruction is a small program stored in an internal memory \\ \hline
    Speed & very fast & a little slower \\ \hline
    Changes & hard (we must change the circuit) & easy (we change the microprogram) \\ \hline
    Found in & RISC processors & CISC processors like the 8086 \\ \hline
\end{tabularx}
\end{center}

%================================================
\section{The instruction set}
%================================================

\begin{tarif}
    The \term{instruction set}{jeu d'instructions}: the list of all the instructions a processor can run.
    Each processor family has its own set.
\end{tarif}

\subsection{The shape of an instruction}
Every machine-language instruction has two parts:

\begin{center}
\begin{tikzpicture}
    \node[draw=brandOrange, fill=brandOrange!20, thick, minimum width=3.5cm, minimum height=1cm, font=\bfseries] (op) at (0,0) {opcode};
    \node[draw=brandTeal, fill=brandTeal!20, thick, minimum width=5cm, minimum height=1cm, font=\bfseries, anchor=west] (oper) at (op.east) {operands};
    \node[below=0.15cm of op, font=\small] {what do we do?};
    \node[below=0.15cm of oper, font=\small] {on what do we work?};
    \node[above=0.2cm of op, font=\codefont] {ADD};
    \node[above=0.2cm of oper, font=\codefont] {AX, 5};
\end{tikzpicture}
\end{center}

\begin{tarif}
    \term{Assembly language}{langage d'assemblage}: we write instructions with short words that humans
    understand (\code{ADD}, \code{MOV}) instead of binary numbers. A program called the
    \term{assembler}{assembleur} turns them into machine language (zeros and ones).
\end{tarif}

\begin{mithal}
    On the 8086, the instruction \code{add al, 5} is written in machine language with the two bytes 04 05, that is
    0000\,0100 0000\,0101: the first byte is the opcode ``add an immediate value to AL'', and the second one is the
    number 5. The programmer writes \code{add al, 5}, and the assembler writes the zeros and ones (we see this in
    detail in Chapter 3).
\end{mithal}

\subsection{How many instructions can a processor have?}
If the opcode uses $k$ bits, there are $2^k$ different codes, so at most $2^k$ instructions.

\begin{khatwa}
    A 6-bit opcode: $2^6 = 64$ possible instructions. An 8-bit opcode: $2^8 = 256$ instructions.
\end{khatwa}

\subsection{Number of addresses in an instruction}
Let us compute $A = B + C$ on four kinds of machines. The fewer operands in an instruction, the shorter the
instruction, but the more instructions we need:

\begin{center}
\small
\begin{tabularx}{\linewidth}{|>{\centering\arraybackslash}p{3.6cm}|X|>{\centering\arraybackslash}p{3cm}|}
    \hline
    \ras{Kind of machine} & \ras{Program} & \ras{Instructions} \\ \hline
    3 addresses & \code{ADD A, B, C} & 1 \\ \hline
    2 addresses & \code{MOV A, B} then \code{ADD A, C} & 2 \\ \hline
    1 address (accumulator) & \code{LOAD B} then \code{ADD C} then \code{STORE A} & 3 \\ \hline
    0 address (stack) & \code{PUSH B}, \code{PUSH C}, \code{ADD}, \code{POP A} & 4 \\ \hline
\end{tabularx}
\end{center}
The 8086 is a \textbf{2-address} machine: \code{add ax, bx} means $\text{AX} \leftarrow \text{AX} + \text{BX}$.

\subsection{Classes of instructions}
\begin{center}
\begin{tabularx}{\linewidth}{|>{\centering\arraybackslash}p{3.5cm}|R|>{\centering\arraybackslash}p{4.2cm}|}
    \hline
    \ras{Class} & \ras{Job} & \ras{Examples (8086)} \\ \hline
    data transfer & copy a value from one place to another & \code{MOV}, \code{PUSH}, \code{POP} \\ \hline
    arithmetic & add, subtract, multiply, divide & \code{ADD}, \code{SUB}, \code{MUL}, \code{INC} \\ \hline
    logic & operations on bits & \code{AND}, \code{OR}, \code{XOR}, \code{SHL} \\ \hline
    flow control & jump and change the order of execution & \code{JMP}, \code{JE}, \code{CALL}, \code{RET} \\ \hline
    input/output & talk with devices & \code{IN}, \code{OUT}, \code{INT} \\ \hline
\end{tabularx}
\end{center}

\subsection{CISC and RISC}
\begin{itemize}
    \item \textbf{CISC} (Complex Instruction Set Computer): many complex instructions; some do a lot of work. Example: Intel x86.
    \item \textbf{RISC} (Reduced Instruction Set Computer): few, simple, fast instructions, all of the same length.
          Example: ARM in phones, RISC-V.
\end{itemize}

%================================================
\section{Addressing modes}
%================================================

\begin{tarif}
    An \term{addressing mode}{mode d'adressage}: the way we say \textbf{where the data is} that the instruction works on.
\end{tarif}

\begin{tashbih}
    You want to give a book to a friend:
    \textbf{(1)} you put the book in his hand (immediate),
    \textbf{(2)} you say: ``the book is on shelf number 12'' (direct),
    \textbf{(3)} you say: ``look at the paper in my pocket, it has the shelf number'' (indirect),
    \textbf{(4)} you say: ``the third shelf after shelf number 10'' (indexed).
\end{tashbih}

\subsection{The basic modes}
\begin{center}
\small
\begin{tabularx}{\linewidth}{|>{\centering\arraybackslash}p{2.4cm}|>{\centering\arraybackslash}p{3.3cm}|R|}
    \hline
    \ras{Mode} & \ras{8086 example} & \ras{Where is the data?} \\ \hline
    immediate & \code{mov ax, 5} & inside the instruction itself \\ \hline
    register & \code{mov bx, ax} & in a register \\ \hline
    direct & \code{mov cx, [val]} & in memory; its address is written in the instruction \\ \hline
    indirect & \code{mov dx, [si]} & in memory; its address is in a register \\ \hline
    indexed & \code{mov ax, [tab+di]} & address = start address + value of an index register \\ \hline
    relative & \code{jmp next} & address = current address + offset (for jumps) \\ \hline
\end{tabularx}
\end{center}

\subsection{When do we use each mode?}
\begin{itemize}
    \item \textbf{Immediate}: for constants (number of loop turns, start value). The fastest mode, because the value
          came with the instruction, so we do not go to memory. But we cannot change it while the program runs.
    \item \textbf{Register}: for fast calculations between values already inside the processor. Very fast too.
    \item \textbf{Direct}: for variables that have a fixed place in memory (like a global variable). One memory read.
    \item \textbf{Indirect}: when we do not know the address when we write the program, but compute it while it runs.
          It is the basis of \textbf{pointers} in C.
    \item \textbf{Indexed}: for arrays. We fix the start address and change the index at each turn of the loop, to visit every element.
    \item \textbf{Relative}: for jumps. Because the distance is relative, the program works even if it is loaded at another place in memory.
\end{itemize}

\begin{khatwa}[: how many times do we go to memory?]
    \begin{center}
    \begin{tabular}{|l|c|}
        \hline \ras{Instruction} & \ras{Memory reads to get the data} \\ \hline
        \code{mov ax, 5} (immediate) & 0 \\ \hline
        \code{mov ax, bx} (register) & 0 \\ \hline
        \code{mov ax, [val]} (direct) & 1 \\ \hline
        \code{mov ax, [si]} (indirect with a register) & 1 \\ \hline
        memory indirect (the address itself is in memory) & 2 \\ \hline
    \end{tabular}
    \end{center}
    Every trip to memory costs time, so programmers try to keep important data in registers.
\end{khatwa}

\subsection{The modes in a real program}
\asm{ch2/modes.asm}

\begin{center}
\begin{tikzpicture}[cellx/.style={draw=brandDark, fill=white, minimum width=2cm, minimum height=0.55cm, font=\codefont\small},
    addr/.style={font=\codefont\small, text=tealDeep}]
    \foreach \a/\v/\n [count=\i from 0] in {0117/1234/val, 0119/000A/tab, 011B/0014/{}, 011D/001E/{}} {
        \node[cellx] (c\i) at (0, -0.55*\i) {\v};
        \node[addr, left=0.15cm of c\i] {\a};
        \node[font=\codefont\small, right=0.15cm of c\i] {\n};
    }
    \node[font=\small\bfseries] at (0,0.55) {memory};
    \draw[sahmO] (-2.4,-0.55) -- (-2.4,-1.65) node[midway, left, font=\scriptsize] {+4};
\end{tikzpicture}
\end{center}

\begin{sharh}
    \s{1}{A program of type .COM: it starts at address \hex{100}.}
    \s{3}{\textbf{Immediate}: the number 5 is written inside the instruction. Result: AX = 0005.}
    \s{4}{\textbf{Register}: copy AX into BX. Result: BX = 0005.}
    \s{5}{\textbf{Direct}: read memory at the address of \code{val} (\hex{0117}). Result: CX = 1234.}
    \s{6}{\code{lea} puts the \textbf{address} of \code{val} (not its value) into SI: SI = 0117.}
    \s{7}{\textbf{Indirect}: read memory at the address held in SI. Result: DX = 1234.}
    \s{8}{DI = 4: each element of \code{tab} is a word (two bytes), so the third element is 4 bytes away.}
    \s{9}{\textbf{Indexed}: address $= \hex{0119} + 4 = \hex{011D}$, which holds 30. Result: AX = 001E.}
    \s{10}{End of the program, back to the system.}
\end{sharh}

\begin{intabih}
    The brackets \code{[ ]} always mean: ``the \textbf{content of memory} at this address''. \code{mov dx, si}
    copies the number \hex{0117} itself, but \code{mov dx, [si]} reads what is stored at address \hex{0117}.
\end{intabih}

%================================================
\section{The steps of an instruction}
%================================================

Every instruction, even a simple one, goes through the same steps. This is called the
\term{instruction cycle}{cycle d'instruction}.

\begin{center}
\begin{tikzpicture}
    \foreach \a/\t/\c in {90/{1. Fetch}/brandTeal, 0/{2. Decode}/brandOrange,
                          -90/{3. Execute}/brandGreen, 180/{4. Write back}/brandPurple} {
        \node[unit=\c, minimum width=2.9cm, minimum height=1.1cm] (n\a) at (\a:2.6) {\t};
    }
    \draw[sahmT] (n90) to[bend left=25] (n0);
    \draw[sahmT] (n0) to[bend left=25] (n-90);
    \draw[sahmT] (n-90) to[bend left=25] (n180);
    \draw[sahmT] (n180) to[bend left=25] (n90);
    \node[font=\small, text width=3cm, align=center] at (0,0) {then again\\with the next instruction};
\end{tikzpicture}
\end{center}

\begin{enumerate}
    \item \textbf{Fetch}: CO is sent to RAM, memory is read, the instruction arrives in RIM, then in RI.
          Then CO increases to point to the next instruction.
    \item \textbf{Decode}: the decoder reads the opcode in RI and finds the operands.
    \item \textbf{Execute}: if we need data from memory we read it, then the UAL computes the result.
    \item \textbf{Write back}: we store the result in a register or in memory, and we update the flags.
\end{enumerate}

\subsection{A complete example, step by step}
We take a simple machine with an accumulator ACC, and we want to compute $z = x + y$. The variable $x$
is at address 20, $y$ at 21, and $z$ at 22.

\begin{center}
\begin{tabular}{|c|l|l|}
    \hline
    \ras{Address} & \ras{Content} & \ras{Meaning} \\ \hline
    10 & \code{LOAD 20}  & $\text{ACC} \leftarrow [20]$ \\ \hline
    11 & \code{ADD 21}   & $\text{ACC} \leftarrow \text{ACC} + [21]$ \\ \hline
    12 & \code{STORE 22} & $[22] \leftarrow \text{ACC}$ \\ \hline
    20 & 7 & $x$ \\ \hline
    21 & 5 & $y$ \\ \hline
    22 & 0 & $z$ (before running) \\ \hline
\end{tabular}
\end{center}

We follow the second instruction \code{ADD 21} step by step (at the start, CO = 11 and ACC = 7):

\begin{center}
\small
\begin{tabularx}{\linewidth}{|>{\centering\arraybackslash}p{1.6cm}|R|c|c|c|c|c|}
    \hline
    \ras{Step} & \ras{What happens?} & \ras{CO} & \ras{RAM} & \ras{RIM} & \ras{RI} & \ras{ACC} \\ \hline
    fetch & \nql{RAM}{CO} & 11 & \cell{orangeDeep}{11} & -- & -- & 7 \\ \cline{2-7}
     & read memory: \nql{RIM}{[11]} & 11 & 11 & \cell{orangeDeep}{\code{ADD 21}} & -- & 7 \\ \cline{2-7}
     & \nql{RI}{RIM}, and \nql{CO}{CO + 1} & \cell{orangeDeep}{12} & 11 & \code{ADD 21} & \cell{orangeDeep}{\code{ADD 21}} & 7 \\ \hline
    decode & the decoder understands: ``add the content of address 21'' & 12 & 11 & \code{ADD 21} & \code{ADD 21} & 7 \\ \hline
    execute & \nql{RAM}{21}, then read: \nql{RIM}{[21]} & 12 & \cell{orangeDeep}{21} & \cell{orangeDeep}{5} & \code{ADD 21} & 7 \\ \cline{2-7}
     & UAL: $7 + 5$ & 12 & 21 & 5 & \code{ADD 21} & 7 \\ \hline
    write back & \nql{ACC}{12}, and update the flags & 12 & 21 & 5 & \code{ADD 21} & \cell{orangeDeep}{\textbf{12}} \\ \hline
\end{tabularx}
\end{center}

\begin{khulasa}
    After this instruction: ACC = 12 and CO = 12, so the next instruction is \code{STORE 22}, which writes 12
    in memory, and $z$ becomes 12.
\end{khulasa}

\begin{hal}
    Modern processors do not wait for one instruction to finish before starting the next. While one instruction
    runs, the next one is decoded and the third one is fetched. This is the \textbf{pipeline}, which we study in Chapter 4.
\end{hal}

%================================================
\section{Exercises}
%================================================

\begin{tamrin}
    Compute with the full adder (bit by bit) $1011 + 0110$ on 4 bits. What is the result, and the value of CF?
\end{tamrin}
\begin{hall}
    \begin{tabular}{|c|c|c|c|c|c|}
        \hline \ras{Bit} & \ras{$A$} & \ras{$B$} & \ras{$C_{in}$} & \ras{$S$} & \ras{$C_{out}$} \\ \hline
        0 & 1 & 0 & 0 & 1 & 0 \\ \hline 1 & 1 & 1 & 0 & 0 & 1 \\ \hline
        2 & 0 & 1 & 1 & 0 & 1 \\ \hline 3 & 1 & 0 & 1 & 0 & 1 \\ \hline
    \end{tabular}\\
    The result is $0001$ and CF = 1: $11 + 6 = 17 = 16 + 1$, and the ``16'' went into CF.
\end{hall}

\begin{tamrin}
    On 8 bits, give AL, CF and ZF after: (a) $255 + 1$ (b) $10 - 10$ (c) $10 - 20$.
\end{tamrin}
\begin{hall}
    (a) $256 - 256 = 0$: AL = 0, CF = 1, ZF = 1.\\
    (b) AL = 0, CF = 0, ZF = 1.\\
    (c) $10 - 20 = -10$, and on 8 bits $256 - 10 = 246 = \hex{F6}$: CF = 1 (borrow), ZF = 0.
\end{hall}

\begin{tamrin}
    Give the addressing mode of each instruction: (a) \code{mov al, [bx]} (b) \code{mov al, 10h} (c) \code{mov al, [10h]} (d) \code{mov al, bl}.
\end{tamrin}
\begin{hall}
    (a) register indirect \quad (b) immediate \quad (c) direct \quad (d) register.
\end{hall}

\begin{tamrin}
    In the example above ($z = x + y$), write the table of steps for the third instruction \code{STORE 22}.
\end{tamrin}
\begin{hall}
    \textbf{Fetch}: \nql{RAM}{12}, \nql{RIM}{STORE 22}, \nql{RI}{RIM}, CO = 13.
    \textbf{Decode}: ``write ACC at address 22''.
    \textbf{Execute and write}: \nql{RAM}{22}, \nql{RIM}{ACC} (12), then a write signal: $[22] = 12$.
\end{hall}
